\bookchapter{Capítulo 62: Desenvolvimento de interface — introdução ao SwiftUI}

\chapterintro{A interface é tudo o que o usuário vê. O SwiftUI é o framework de interface declarativa lançado pela Apple em 2019 — você apenas “declara” como a interface deve ficar, e o sistema cuida sozinho da renderização, do modo escuro e da adaptação às várias telas. Em 2026, o SwiftUI já é o padrão absoluto para projetos novos. Neste capítulo, vamos aprender a “pintar”.}

\booksection{A ideia do SwiftUI: declarativo, não imperativo}

O desenvolvimento de interface antigo é “imperativo”: você diz ao sistema “crie um botão > coloque em (20,30) > defina a cor”, dando ordens passo a passo. O SwiftUI é “declarativo”: você simplesmente descreve “na tela há um botão vermelho escrito ‘Toque aqui’”, e o sistema se vira sozinho para desenhar e redesenhar quando for preciso.

\begin{lstlisting}[language=Swift]
import SwiftUI

struct ContentView: View {
    var body: some View {      // body: descreve "como a tela fica"
        VStack {               // empilhado na vertical
            Text("Olá, mundo!")
                .font(.largeTitle)      // modificador: título grande
            Button("Toque aqui") {
                print("clicado")
            }
            .buttonStyle(.borderedProminent)  // botão azul sólido
        }
        .padding()
    }
}
\end{lstlisting}

Ponto-chave: uma view é uma struct que adota o protocolo View, a propriedade body descreve a interface e os .modificadores() se encadeiam como “maquiagem”. Esse é todo o alicerce do SwiftUI.

\booksection{Views básicas: texto, botão, campo de texto e imagem}

\begin{center}
\begin{tabularx}{\linewidth}{|X|X|X|}
\hline
\textbf{View} & \textbf{Para que serve} & \textbf{Exemplo} \tabularnewline
\hline
Text & Mostrar texto & Text("Olá") \tabularnewline
Button & Botão & Button("Toque aqui") \{ ação \} \tabularnewline
TextField & Campo de texto & TextField("Dica", text: \$texto) \tabularnewline
Image & Mostrar imagem & \texttt{Image(systemName: "star.\allowbreak fill")} \tabularnewline
\hline
\end{tabularx}
\end{center}

\begin{center}\small\color{gray}As quatro views básicas do SwiftUI\end{center}

Repare no Image(systemName:) — a Apple inclui a biblioteca de ícones SF Symbols (milhares de ícones gratuitos): um nome como “star.fill” já exibe uma estrela preenchida na hora. É um benefício escondido do desenvolvimento iOS: para montar interfaces, quase não é preciso procurar ícones.

\booksection{O trio de layout: VStack, HStack e ZStack}

\begin{itemize}
  \item VStack: organiza na vertical (como uma fila do bandejão);
  \item HStack: organiza na horizontal (como gente sentada lado a lado);
  \item ZStack: empilha (como cartas de baralho umas sobre as outras, ótimo para selos e sobreposições).
\end{itemize}

\begin{lstlisting}[language=Swift]
VStack(alignment: .leading, spacing: 12) {
    HStack {
        Text("Conta")
        Spacer()                 // espaço elástico: empurra a direita até o fim
        Text("-25.5 yuan").foregroundColor(.red)
    }
    Divider()                    // divisor
    Text("Obs.: almoço")
}
.padding(16)                     // espaço nas bordas
.background(Color(.systemGray6)) // fundo cinza claro
.cornerRadius(12)                // cantos arredondados

// A ordem dos modificadores importa: padding primeiro, depois background — o fundo não inclui o espaço
\end{lstlisting}

\begin{tcbwarn}{Atenção: a ordem dos modificadores importa}
padding(16).background(cinza) e background(cinza).padding(16) dão resultados diferentes — o primeiro “deixa o espaço e depois pinta” (o fundo não inclui a margem), o segundo “pinta e depois deixa o espaço” (o fundo inclui a margem). Daí vêm as “bordas brancas estranhas” que tanto atormentam iniciantes. Na dúvida, inverta a ordem: a pré-visualização responde na hora.

\end{tcbwarn}

\booksection{Estado e binding: dando vida à interface}

Interface estática não tem graça. A regra do SwiftUI: a interface reflete fielmente o estado, e quando o estado muda, ela se redesenha sozinha.

\begin{lstlisting}[language=Swift]
struct CounterView: View {
    @State private var count = 0   // @State: se esse valor muda, redesenha

    var body: some View {
        VStack {
            Text("Você tocou \(count) vezes")
            Button("Toque +1") {
                count += 1          // muda o estado; a tela atualiza sozinha
            }
        }
    }
}
\end{lstlisting}

O campo de texto usa “binding bidirecional” — o símbolo \$ entrega junto o valor e o poder de alterá-lo:

\begin{lstlisting}[language=Swift]
struct InputView: View {
    @State private var text = ""

    var body: some View {
        TextField("Digite algo", text: $text)  // $ vínculo bidirecional
            .textFieldStyle(.roundedBorder)
        Text("Você digitou: \(text)")   // atualiza em tempo real
    }
}
\end{lstlisting}

Ação do usuário

O estado muda (count += 1)

O SwiftUI percebe a mudança

A tela redesenha sozinha (body recalculado)

\begin{center}\small\color{gray}Figura 3: o fluxo de dados unidirecional “estado interface” do SwiftUI\end{center}

\booksection{Prévia, tema e animação: os “cheats” do SwiftUI}

\begin{itemize}
  \item Macro \#Preview: adicione \#Preview \{ ContentView() \} no fim do arquivo e o canvas à direita renderiza em tempo real — editar vendo o resultado na hora;
  \item Modo escuro de graça: use cores do sistema (Color(.systemGray6), .primary) em vez de cores fixas e o modo escuro se adapta sozinho, sem uma linha a mais;
  \item Animação em uma linha: withAnimation \{ count += 1 \} ou .animation(.default, value: count) na view — transições sedosas.
\end{itemize}

\begin{lstlisting}[language=Swift]
struct AnimatedHeartView: View {
    @State private var liked = false

    var body: some View {
        Button {
            withAnimation(.spring(duration: 0.4)) {
                liked.toggle()
            }
        } label: {
            Image(systemName: liked ? "heart.fill" : "heart")
                .font(.system(size: 48))
                .foregroundColor(liked ? .red : .gray)
                .scaleEffect(liked ? 1.3 : 1.0)
        }
    }
}
\end{lstlisting}

\begin{tcbinfo}{Dica: primeiro a função, depois a animação}
O defeito clássico do iniciante é “ainda não ter acertado a interface e já ficar remoendo a animação”. Faça a função funcionar primeiro e deixe os efeitos para o final — animação é enfeite no bolo pronto; função é remédio na hora da dor. É também o nosso princípio de sempre: simplicidade acima de tudo.

\end{tcbinfo}

\begin{tcbpractice}{Mãos à obra: crie um “botão do humor”}
\booksection{Consulta rápida dos modificadores mais usados: com esta página, a interface se resolve}

\begin{center}
\begin{tabularx}{\linewidth}{|X|X|}
\hline
\textbf{O que você quer fazer} & \textbf{Modificador} \tabularnewline
\hline
Espaçamento & \texttt{.\allowbreak padding(16)} \tabularnewline
Cantos arredondados & \texttt{.\allowbreak cornerRadius(12)} \tabularnewline
Cor de fundo & \texttt{.\allowbreak background(Color(.\allowbreak systemGray6))} \tabularnewline
Fonte e tamanho & \texttt{.\allowbreak font(.\allowbreak title.\allowbreak bold())}/\texttt{.\allowbreak font(.\allowbreak system(size: 20))} \tabularnewline
Cor de primeiro plano (texto/ícone) & \texttt{.\allowbreak foregroundColor(.\allowbreak red)} \tabularnewline
Ocupar toda a largura & \texttt{.\allowbreak frame(maxWidth: .\allowbreak infinity)} \tabularnewline
Tamanho fixo & \texttt{.\allowbreak frame(width: 100, height: 44)} \tabularnewline
Alinhamento & \texttt{VStack(alignment: .\allowbreak leading)} \tabularnewline
Transparência & \texttt{.\allowbreak opacity(0.\allowbreak 6)} \tabularnewline
Sombra & \texttt{.\allowbreak shadow(radius: 5)} \tabularnewline
Gesto de toque & \texttt{.\allowbreak onTapGesture \{ .\allowbreak .\allowbreak .\allowbreak  \}} \tabularnewline
Tipo de teclado & \texttt{.\allowbreak keyboardType(.\allowbreak decimalPad)} \tabularnewline
\hline
\end{tabularx}
\end{center}

\begin{center}\small\color{gray}Consulta rápida dos modificadores mais usados do SwiftUI\end{center}

Se não decorar tudo, sem problema — esta página existe para consulta. Na prática, os quatro primeiros são os mais usados: padding, cantos arredondados, fundo e fonte.

Transforme o projeto modelo em um alternador de humor: um botão que, a cada toque, cicla entre feliz triste bravo, com o texto e a cor de fundo mudando junto. Dica: guarde os três pares (texto + cor) em um array, o índice em um @State e, ao tocar, use (index + 1) \% 3. Mande para os amigos testarem e colha a primeira leva de “uau, você faz apps!”.

\end{tcbpractice}

\begin{tcbsummary}{Resumo do capítulo}
\begin{itemize}
  \item O SwiftUI é declarativo: você descreve “como deve ficar”, e o sistema desenha e redesenha
  \item View = struct que adota o protocolo View; o body descreve a interface; os modificadores se encadeiam como “maquiagem”
  \item As quatro views básicas: Text / Button / TextField / Image; a biblioteca SF Symbols dá ícones de graça e de montão
  \item O trio de layout: VStack empilha na vertical, HStack na horizontal, ZStack em camadas; o Spacer é o espaço elástico
  \item @State dá vida à interface; o \$ faz o binding bidirecional que alimenta o campo de texto
  \item \#Preview mostra o resultado na hora; as cores do sistema se adaptam ao modo escuro sozinhas; withAnimation dá o efeito em uma linha
\end{itemize}

\end{tcbsummary}
